%! AP Statistics — Unit 7, Lesson 7.1 — Homework
\documentclass[10pt]{article}
\usepackage{apstats-article}
\usepackage{apstats-boxes}
\usepackage{grading}
\gradingassignment{apstats-u07-l01-hw}
\gradingcourse{advanced_statistics}

\begin{document}

\pageheader{Unit 7, Lesson 7.1}{Homework}
\namedateperiod

\vspace{0.1in}
\textbf{Directions:} For each scenario, complete all parts. Show your reasoning clearly.

\begin{enumerate}

  \item \textbf{Airport Security.} A transportation agency tests whether the mean wait
        time at a security checkpoint exceeds 15 minutes
        ($H_0{:}\ \mu = 15$ min, $H_a{:}\ \mu > 15$ min, $\alpha = 0.05$).

        \begin{enumerate}[label=\alph*.]
          \item Describe a Type I error in context.
                \gradedhide{\writeline\\\gradedhere{\gradedblank[id=hw1a, type=rubric, points=2]{\linewidth+6pt}{\gradedlinesheight{1}}}\writelines{1}}
          \item Describe a Type II error in context.
                \gradedhide{\writeline\\\gradedhere{\gradedblank[id=hw1b, type=rubric, points=2]{\linewidth+6pt}{\gradedlinesheight{1}}}\writelines{1}}
          \item Which error is more consequential for travelers? Justify your answer.
                \gradedhide{\writeline\\\gradedhere{\gradedblank[id=hw1c, type=rubric, points=2]{\linewidth+6pt}{\gradedlinesheight{1}}}\writelines{1}}
        \end{enumerate}

  \item \textbf{Online Learning Platform.} A school district tests whether the mean
        score on a state exam differs for students who used an online learning platform
        compared to the district average of 72 points
        ($H_0{:}\ \mu = 72$, $H_a{:}\ \mu \ne 72$, $\alpha = 0.05$).

        \begin{enumerate}[label=\alph*.]
          \item Describe a Type I error in context.
                \gradedhide{\writeline\\\gradedhere{\gradedblank[id=hw2a, type=rubric, points=2]{\linewidth+6pt}{\gradedlinesheight{1}}}\writelines{1}}
          \item Describe a Type II error in context.
                \gradedhide{\writeline\\\gradedhere{\gradedblank[id=hw2b, type=rubric, points=2]{\linewidth+6pt}{\gradedlinesheight{1}}}\writelines{1}}
          \item The district administrator wants to lower $\alpha$ to 0.01.
                Explain the trade-off this creates.
                \gradedhide{\writeline\\\gradedhere{\gradedblank[id=hw2c, type=rubric, points=2]{\linewidth+6pt}{\gradedlinesheight{1}}}\writelines{1}}
        \end{enumerate}

  \item \textbf{Environmental Testing.} A water treatment facility tests whether the
        mean lead concentration in tap water is below a safe threshold of 10 ppb
        ($H_0{:}\ \mu = 10$ ppb, $H_a{:}\ \mu < 10$ ppb, $\alpha = 0.05$).

        \begin{enumerate}[label=\alph*.]
          \item Describe both error types in context.
                \gradedhide{\writeline\\\gradedhere{\gradedblank[id=hw3a, type=rubric, points=2]{\linewidth+6pt}{\gradedlinesheight{2}}}\writelines{2}}
          \item In this context, which error type is more dangerous? Explain.
                \gradedhide{\writeline\\\gradedhere{\gradedblank[id=hw3b, type=rubric, points=2]{\linewidth+6pt}{\gradedlinesheight{1}}}\writelines{1}}
        \end{enumerate}

\end{enumerate}

\begin{reflectionbox}
\textbf{Reflection:} Explain in your own words why it is impossible to simultaneously
minimize both the Type I and Type II error rates for a fixed sample size.

\writelines{3}
\end{reflectionbox}

\end{document}
